%!TEX root = ../LLMSeminar.tex
% ──────────────────────────────────────────────────────────────
%  Section 9 — Calibrating Expectations: The Three-Tier Exercise
% ──────────────────────────────────────────────────────────────

\section{Calibrating expectations: the three-tier exercise}%
\label{sec:three-tier}

\subsection{Why a personal benchmark}

The popular discourse around \emph{artificial general
intelligence} is, to put it gently, broken.  No participant
agrees on what counts as ``general,'' and no participant
agrees on what counts as ``intelligence.''  The result is a
debate impossible to ground.

In a working seminar this is unhelpful.  We need an internal
sense of what a current model can and cannot do, sharp
enough that we are not surprised by genuine progress and
sharp enough that we are not impressed by mere theatre.  The
calibration cannot come from headlines; it has to come from
experiments \emph{you} run, on tasks \emph{you} care about.

\begin{keyresult}[The three-tier exercise]\label{kr:three-tier}
  Each week, choose three tasks across a spectrum of difficulty:
  \begin{enumerate}[(T1)]
    \item Something you would be \emph{mildly impressed} to
      see an AI do.  ``Yes, I sort of expected that to work.''
    \item Something you would be \emph{well impressed} by.
      ``Now that I see it, this is genuinely useful.''
    \item Something where, if it worked, \emph{the world has
      changed}.  ``If this works, our intuitions need to be
      rebuilt.''
  \end{enumerate}
  Try them.  Record the outcome.  Repeat the exercise on a
  different week with the same three tasks.  This is your
  personal capability index.
\end{keyresult}

\begin{intuition}[Why three tiers, not one]
  A single test gives a binary outcome and tells you almost
  nothing about the slope.  The interesting question is
  \emph{where the frontier sits}: which kinds of problems are
  solid, which are at the edge, and which look impossible
  today.  A spread of three tests gives a coarse but useful
  derivative.  Re-run them quarterly and you can watch the
  frontier move.
\end{intuition}

\subsection{What the room produced}

The audience, working live, contributed the following examples.
They are recorded here both as data on contemporary capability
and as suggestions for the reader's own experiments.

\subsubsection{T3: ``Earn enough to pay for itself''}

A wonderfully meta task.  Give an agent the prompt:

\begin{lstlisting}
You are autonomous.  Acquire money on the open Internet.
You may not start with funds.  Stop when you have generated
enough revenue to pay for your own API usage.
\end{lstlisting}

The implicit problem set is enormous: open a bank or wallet
account, accept payment, route through some legal economic
activity, and break even.  Each sub-problem is itself
non-trivial; the integral is currently somewhere between
``hard'' and ``impossible.''

\begin{intuition}[A moving target]
  This task is moving.  Every effective method of low-friction
  income on the Internet eventually saturates as people---and
  agents---exploit it.  An agent that succeeds today may fail
  next year, and one that fails today may succeed in two.
  This is part of what makes it a useful T3 benchmark: it tracks
  not just the agent but the agent's environment.
\end{intuition}

\subsubsection{T1: handwritten lecture notes $\to$ \LaTeX}

A confirmed positive result.  Photographs of handwritten
notes (mathematics included) were uploaded to a frontier
chat model and transcribed to typeset \LaTeX{} faithfully---
equations, indices, and structure preserved.  The remaining
human work was the \emph{exposition}: the connective text
between equations.  This is doable today, and it is exactly
how I produced the working version of my general-relativity
lecture script: handwritten notes, plus pandemic-era video
lectures (transcribed by Whisper~\cite{radford2023whisper}),
plus an LLM that filled in the exposition between equations
using the spoken commentary as ground truth.

\begin{remark}
  This workflow is recommended for any researcher who has
  written legible handwritten notes that they would prefer
  not to retype.  The bottleneck is that, on its own, an
  image-to-\LaTeX{} pass produces only equations; the
  surrounding text has to come from somewhere.  An audio
  recording of the original lecture supplies it.
\end{remark}

\subsubsection{T2: a Signal plug-in for in-line \LaTeX}

A user wants to type \verb|$E = mc^2$| in Signal and see it
rendered.  Signal is open in code but closed in server
access; plug-ins are not officially supported.  As a
contained task this is solidly in T2 territory---non-trivial
research and code, but no clear blocker---unless one were
to require fully native deployment, in which case the
licensing and server architecture push it towards T3.

\subsubsection{T2: AR poster translator}

Conference poster sessions are walls of text.  A camera
feed that overlaid clean, summarised text in real time
(in the spirit of Google Lens, but for academic posters)
would be transformative for any large meeting.  Underlying
components---OCR, translation, summarisation---are
solved; the integration into a phone-based experience is
engineering.  T2.

\subsubsection{T2: open problems on the arXiv, ranked}

Build an index of \emph{actually open} problems by mining
the conclusions sections of arXiv papers, ranked by how
many times each problem is referenced, and cross-referenced
against subsequent papers that claim to solve them.
Doable in principle today; the limiting factor is token
budget rather than capability.

\subsubsection{T2: reschedule a course via Stud-IP's API}

A live demonstration: Claude was asked to investigate
the Stud-IP learning-management system and figure out how
to change the day on which a course meets.  It scraped the
documentation, identified the JSON\:API endpoints, and
produced a working shell command---halting just short of
executing it because permissions had not been granted.
This is precisely the boring-but-genuinely-useful work that
fills an academic week.

\subsubsection{T1--T2: 3D model from photographs}

Photogrammetry is well-developed; with sufficiently many
images, recovering a 3D mesh is essentially solved.  The
interesting frontier is doing it from very few images by
relying on object priors---knowing what kind of thing is in
the picture.  Public products already do this for
landmarks (Google Maps' 3D buildings); doing it for arbitrary
objects with $3$--$5$ photos is at the research frontier
but moving fast.

\subsubsection{T2--T3: full-symphony automatic music transcription}

The audience asked: take an audio recording of a Mahler
symphony and produce the full sheet music for every
instrument.  When asked, the agent itself surveyed the state
of the art and reported, accurately:

\begin{itemize}
  \item Solo piano transcription is largely solved
    \cite{benetos2018amt,hawthorne2018magenta}.
  \item Small ensembles are partially solved with
    non-trivial error rates.
  \item Full orchestral source separation is \emph{not}
    solved: state-of-the-art models confuse violins with
    violas, fail to separate desks within a section, and
    degrade rapidly past about six pitches.
\end{itemize}

This is the genuinely informative outcome.  Not because the
agent solved the problem (it cannot), but because it
correctly identified the frontier and articulated why the
problem is hard.  The technical bottleneck here lives outside
the LLM: it is in the underlying audio models and the
training data available for them.

\subsubsection{T3: physical inference from a single audio recording}

A genuinely surprising positive result reported by an
audience member.  A short ($\sim 10$~s) recording of
\emph{ambient} room sound was given to a frontier coding
agent with the instruction:

\begin{lstlisting}
Use every signal-processing and acoustic-inference tool
available to extract whatever you can about the
environment in which this was recorded.
\end{lstlisting}

The agent inferred:

\begin{itemize}
  \item Approximate room dimensions (length, width, height)
    from resonant modes;
  \item Floor and ceiling materials from spectral colouration
    and absorption profile;
  \item The presence of large furniture from
    early-reflection patterns;
  \item Occupancy and ``furnished-ness'' from the room's
    overall absorption coefficient ($T_{60}$);
  \item A geographical hint (parts of Europe) from the bird
    species audible in the background, and a finer one from
    a public-address announcement also picked up;
  \item With small motion (footsteps), Doppler-style
    inversion of the time-varying response to triangulate
    fixed objects in the room.
\end{itemize}

\begin{intuition}[The physics is real]
  This is, at heart, an inverse problem in acoustics:
  given the impulse response of a cavity, recover the
  cavity's geometry and surface impedances.  The literature
  on \emph{room geometry estimation from impulse
  responses}~\cite{krijnen2014roomgeom,antonacci2012geometric}
  is decades old; modern acoustic emulators
  \cite{scheibler2018pyroomacoustics} make forward modelling
  cheap.  What is new is that an off-the-shelf coding agent
  can autonomously assemble the pipeline.  The scientific
  content is honest physics.
\end{intuition}

\begin{remark}[On disclosure]
  Several of these inferences are dual-use.  Discretion
  is appropriate when discussing them outside the seminar
  room.
\end{remark}

\subsection{The shape of the lesson}

\begin{keyresult}[What the exercise reveals]
  Across the dozen problems collected, the rough pattern was:
  \begin{itemize}
    \item \textbf{T1 mostly solved.}  Tasks that combine
      established perception (OCR, ASR, translation) with
      structured output (\LaTeX, JSON) work reliably.
    \item \textbf{T2 is contingent.}  Many tasks land here:
      doable, but with serious engineering or permission
      barriers.  Whether the agent ``does it'' depends on
      access.
    \item \textbf{T3 is the interesting tier.}  Some T3 tasks
      proved surprisingly tractable (acoustic inference);
      others, where verification is the hard part (full
      symphony AMT, autonomous fund-raising), remain
      genuinely beyond reach.
  \end{itemize}
\end{keyresult}

\begin{audience}[Audience consensus]
  We will not be able to define ``intelligence.''  We
  \emph{can} build intuition by running the experiments and
  watching where they pass and fail.  This is the empirical
  posture the seminar adopts going forward.
\end{audience}
